% Copyright (c) 2026 Benjamin Stanley Frohman.
% Licensed under the Apache License, Version 2.0.
% REVIEW DRAFT --- not a DOI assignment, upload, or publication.
\documentclass[11pt,a4paper]{article}

\usepackage[margin=1.05in]{geometry}
\usepackage{amsmath,amssymb,amsthm,mathtools}
\usepackage{microtype}
\usepackage{enumitem}
\usepackage{booktabs}
\usepackage{longtable}
\usepackage{hyperref}
\usepackage[T1]{fontenc}
\usepackage{lmodern}
\usepackage{setspace}
\usepackage{xcolor}
\usepackage{xurl}
\Urlmuskip=0mu plus 3mu
\setstretch{1.08}

\hypersetup{colorlinks=true,linkcolor=black,citecolor=black,urlcolor=blue}

\newcommand{\Q}{\mathbb{Q}}
\newcommand{\C}{\mathbb{C}}
\newcommand{\Z}{\mathbb{Z}}
\newcommand{\lean}[1]{\texttt{\detokenize{#1}}}

\newtheorem{theorem}{Theorem}
\newtheorem{lemma}[theorem]{Lemma}
\newtheorem{proposition}[theorem]{Proposition}
\theoremstyle{definition}
\newtheorem{definition}[theorem]{Definition}
\newtheorem{remark}[theorem]{Remark}

\title{\textbf{Polynomial, Determinant, and Formal-Derivative Identities\\
for an Explicit Sextic Section}\\[0.35em]
\large Lean Stage~1 Algebraic Certificates\\[0.4em]
\normalsize\textcolor{red}{\textbf{REVIEW DRAFT}} --- not peer-reviewed; no DOI assigned; not uploaded}
\author{Frohman, Benjamin Stanley\\
\small ORCID: \href{https://orcid.org/0009-0006-7068-3718}{0009-0006-7068-3718}\\
\small Email: \texttt{frohmanbenjamin@gmail.com}\\
\small Independent Researcher}
\date{6~November~2026 (America/Chicago)\\
\small Source pin: \texttt{e207394a9ed4255ad804860b8f69fde7e940dd54}\\
\small Repository: \url{https://github.com/BenFrohman/HODGE}\\
\small Pull request: \url{https://github.com/BenFrohman/HODGE/pull/36}}

\begin{document}
\maketitle

\begin{center}
\fbox{\parbox{0.92\textwidth}{\centering\small
\textbf{Status.} This manuscript is a \textbf{REVIEW DRAFT} prepared for later
integration with Zenodo packaging. It invents no DOI, performs no upload or
publication, and does not claim peer review, novelty, or a Clay/Hodge close.
License of the cited Lean sources remains Apache-2.0; this note does not
relicense them. Author metadata uses the exact form
\texttt{Frohman, Benjamin Stanley}.}}
\end{center}

\begin{abstract}
We record seven kernel-checked commutative-algebra lemmas formalized in Lean~4
in the module \lean{Hodge.ExplicitResidueAlgebra}. The lemmas concern an
exact two-pair decomposition, a binary split factorization, two $4\times4$
Jacobian determinant factorizations, nonvanishing of every root of an explicit
degree-$24$ section modulus in characteristic zero, nonvanishing of one
leading Wronskian coefficient at every such root, and an identity expressing
the binary Wronskian of two explicit binary forms via formal multivariate
polynomial derivatives. The statements are mapped to Lean declarations below. They do \emph{not} construct Betti or Hodge
cohomology, cycle classes, geometric ideals, or a residue-cycle theorem.
Toolchain pins, a theorem-to-declaration table, reproduction commands, and
precise formalization boundaries are included. Companion Stage~2
nonvanishing certificates appear in a separate note.
\end{abstract}

\tableofcontents

\section{Scope and non-claims}

Let $X\subset\mathbb{P}^5$ denote the complex projective hypersurface cut out
by the homogeneous sextic
\begin{equation}
\sum_{j=0}^{2}\bigl(x_j^5\,x_{j+3}+x_{j+3}^6\bigr)=0
\end{equation}
{\sloppy
(indices understood as in the repository's explicit-curve encoding).
The classical manuscript and Python residue certificate in the HODGE
repository argue that certain complete intersection surfaces on~$X$
project nontrivially onto specified Griffiths Jacobian
characters~\cite{HODGERepo}. The present note formalizes only a bounded
algebraic fragment of that route.
}

\paragraph{What this note proves.}
{\sloppy
Exactly the seven Lean theorems listed in
Section~\ref{sec:mapping}, as stated in
\path{Hodge/ExplicitResidueAlgebra.lean} at mathematical
source pin
\path{e207394a9ed4255ad804860b8f69fde7e940dd54}.
}

\paragraph{What this note does not prove.}
\begin{itemize}[leftmargin=1.4em,itemsep=0.25em]
\item The Hodge conjecture for $X$, in any rational or integral form, and
      not the general Hodge conjecture.
\item Geometric cycle classes, Betti cohomology, Hodge structures, or a
      Lean realization of Villaflor Loyola's residue-cycle
      theorem~\cite{Villaflor2021}.
\item An actual Jacobian-quotient vector-space basis, reducedness or
      saturation of geometric ideals $(Q,C)$ or $(L,Q,C)$, or Fourier/Galois
      orbit spanning of rational Hodge classes.
\item The Stage~2 simultaneous nonvanishing of all eight Wronskian/U
      coefficients and the discriminant $b^2-4c$ (treated in the companion
      note), and the tenth Python H/G-resultant certificate (still not
      Lean-proved).
\end{itemize}
These boundaries follow the repository records
\path{docs/FULL_FORMALIZATION_STAGE1.md} and
\lean{docs/FULL_FORMALIZATION_STAGE2.md}~\cite{HODGERepo}.

\section{Notation and standing hypotheses}

\begin{definition}[Section binary forms]
For a commutative ring $R$ and coefficients $b,c,d,e,g\in R$, write
\begin{align}
H(b,c;X,Y) &= X^2 + b\,XY + c\,Y^2,\\
G(d,e,g;X,Y) &= X^3 + d\,X^2Y + e\,XY^2 + g\,Y^3,
\end{align}
as in \lean{Hodge.ExplicitCurve.H} and \lean{Hodge.ExplicitCurve.G}.
\end{definition}

\begin{definition}[Section modulus]
Over any field $K$ of characteristic zero, the scalar section modulus is the
degree-$24$ monic polynomial identity
\begin{equation}\label{eq:modulus}
k^{24}+66081312000\,k^{18}-4811512860000\,k^{12}
+1167566400000\,k^{6}+583200000=0.
\end{equation}
The Lean definitions \lean{ExplicitCurve.sb}, \lean{sc}, \lean{sd},
\lean{se}, \lean{sg} evaluate five explicit rational polynomials in $k$
(recorded in \lean{Hodge/ExplicitCurve.lean}). Stage~1 uses the already
checked relation
\begin{equation}
2\cdot\mathrm{sd}(k)-3\cdot\mathrm{sb}(k)-k=0
\end{equation}
(\lean{ExplicitCurve.rel_h5}) whenever~\eqref{eq:modulus} holds.
\end{definition}

\begin{remark}
The docstring of \lean{ExplicitCurve.m} describes $m$ as irreducible over
$\Q$. That irreducibility claim is classical/Python context and is
\emph{not} among the seven Lean theorems of this note. Stage~2 proves only
$\mathrm{degree}(m)=24$ in Lean.
\end{remark}

\section{Main results}

Throughout, $R$ is a commutative ring and $K$ is a field of characteristic
zero (with a $\Q$-algebra structure when section coefficients appear).

\begin{theorem}[Two-pair decomposition]\label{thm:two-pair}
For $h_1,g_1,h_2,g_2,i\in R$ with $i^2=-1$,
\begin{align*}
&(h_2+h_1)\bigl(-(h_2^2-h_2 h_1+h_1^2)\bigr)+(g_2+i g_1)(g_2-i g_1)\\
&\qquad=(g_1^2-h_1^3)+(g_2^2-h_2^3).
\end{align*}
\end{theorem}
\begin{proof}[Lean proof outline]
\lean{linear_combination} against $-(g_1^2)\,(i^2+1)$. Declaration:
\path{Hodge.ExplicitResidueAlgebra.two_pair_decomposition}.
\end{proof}

\begin{theorem}[Split-pair factorization]\label{thm:split}
For $u,v\in R$,
\begin{equation}
(u+v)\cdot\bigl(v(u^4-u^3 v+u^2 v^2-u v^3+v^4)\bigr)=u^5 v+v^6.
\end{equation}
\end{theorem}
\begin{proof}[Lean proof outline]
Immediate ring normalization (\lean{ring}).
Declaration: \path{split_pair_decomposition}.
\end{proof}

\begin{theorem}[Outer Jacobian determinant]\label{thm:outer}
For $h_1,h_2,i\in R$,
\begin{equation}
\det
\begin{pmatrix}
1 & 0 & 1 & 0\\
h_2-2h_1 & 0 & h_1-2h_2 & 0\\
0 & i & 0 & 1\\
0 & -i & 0 & 1
\end{pmatrix}
=6i(h_2-h_1).
\end{equation}
\end{theorem}
\begin{proof}[Lean proof outline]
Expand via \path{Matrix.det_succ_row_zero} and finish by \lean{ring}.
Declaration: \lean{outer_jacobian_det}. No square-root hypothesis is used.
\end{proof}

\begin{theorem}[Paired Jacobian determinant]\label{thm:paired}
For $h_1,h_2,i\in R$ and eight first partials
\[
h_{1x},\; h_{1y},\; g_{1x},\; g_{1y},\; h_{2x},\; h_{2y},\; g_{2x},\; g_{2y}\in R,
\]
\begin{align*}
&\det
\begin{pmatrix}
h_{1x} & h_{1y} & h_{2x} & h_{2y}\\
(h_2-2h_1)h_{1x} & (h_2-2h_1)h_{1y} & (h_1-2h_2)h_{2x} & (h_1-2h_2)h_{2y}\\
i g_{1x} & i g_{1y} & g_{2x} & g_{2y}\\
-i g_{1x} & -i g_{1y} & g_{2x} & g_{2y}
\end{pmatrix}\\
&\qquad=6i(h_2-h_1)\,(h_{1x}g_{1y}-h_{1y}g_{1x})\,(h_{2x}g_{2y}-h_{2y}g_{2x}).
\end{align*}
\end{theorem}
\begin{proof}[Lean proof outline]
Same determinant expansion and ring normalization as
Theorem~\ref{thm:outer}. Declaration: \lean{paired_jacobian_det}. Linking
the eight symbols to actual polynomial derivatives is a separate obligation
and is not discharged by this identity alone.
\end{proof}

\begin{theorem}[Section roots are nonzero]\label{thm:rootnz}
If $k\in K$ satisfies~\eqref{eq:modulus}, then $k\neq 0$.
\end{theorem}
\begin{proof}[Lean proof outline]
Substitute $k=0$ into~\eqref{eq:modulus} and contradict by
\lean{norm_num}. Declaration: \lean{section_root_ne_zero}.
\end{proof}

\begin{theorem}[Leading $W$ coefficient]\label{thm:leadingW}
Assume moreover that $K$ carries a $\Q$-algebra structure. If $k\in K$ satisfies~\eqref{eq:modulus},
then
\begin{equation}
2\cdot\mathrm{sd}(k)-3\cdot\mathrm{sb}(k)\neq 0.
\end{equation}
\end{theorem}
\begin{proof}[Lean proof outline]
From \lean{ExplicitCurve.rel_h5},
$2\cdot\mathrm{sd}(k)-3\cdot\mathrm{sb}(k)=k$; apply
Theorem~\ref{thm:rootnz}. Declaration: \lean{leading_W_ne_zero}.
This is one of eight coefficient nonvanishings required by the classical
residue route; the remaining seven and the discriminant are Stage~2.
\end{proof}

\begin{theorem}[Binary Wronskian via formal derivatives]\label{thm:wronskian}
For $b,c,d,e,g\in R$, let $x=X_0$, $y=X_1$ in $R[X_0,X_1]$, and set
\begin{align*}
H_p &= x^2 + C(b)\,xy + C(c)\,y^2,\\
G_p &= x^3 + C(d)\,x^2 y + C(e)\,x y^2 + C(g)\,y^3.
\end{align*}
Then the formal partial Wronskian expands as
\begin{align*}
\partial_0 H_p\cdot\partial_1 G_p-\partial_1 H_p\cdot\partial_0 G_p
&= C(2d-3b)\,x^3 + C(4e-bd-6c)\,x^2 y\\
&\qquad + C(6g+be-4cd)\,xy^2 + C(3bg-2ce)\,y^3.
\end{align*}
\end{theorem}
\begin{proof}[Lean proof outline]
Unfold definitions, apply \lean{simp} for the multivariate derivative API,
and finish by \lean{ring}. Declaration: \lean{binary_wronskian}. The
coefficient formulas are exactly the scalar Stage~2 quantities $w_3,w_2,w_1,w_0$.
\end{proof}

\section{Theorem--declaration mapping}
\label{sec:mapping}

\begin{center}
\begin{longtable}{@{}llp{5.6cm}@{}}
\toprule
\# & Lean declaration & Mathematical content \\
\midrule
1 & \lean{two_pair_decomposition} & Thm.~\ref{thm:two-pair}: two-pair sextic identity \\
2 & \lean{split_pair_decomposition} & Thm.~\ref{thm:split}: $u^5v+v^6$ factorization \\
3 & \lean{outer_jacobian_det} & Thm.~\ref{thm:outer}: outer $4\times4$ factor \\
4 & \lean{paired_jacobian_det} & Thm.~\ref{thm:paired}: paired $4\times4$ factor \\
5 & \lean{section_root_ne_zero} & Thm.~\ref{thm:rootnz}: $k\neq0$ \\
6 & \lean{leading_W_ne_zero} & Thm.~\ref{thm:leadingW}: $2d-3b\neq0$ \\
7 & \lean{binary_wronskian} & Thm.~\ref{thm:wronskian}: formal Wronskian \\
\bottomrule
\end{longtable}
\end{center}

\noindent
Namespace: \lean{Hodge.ExplicitResidueAlgebra}. Source file:
\path{Hodge/ExplicitResidueAlgebra.lean}. No generated helper declarations
from Stage~2 are split into this paper.

\section{Axiom audit and verification evidence}
{\sloppy

A Stage~1 seven-declaration \lean{\#print axioms} audit
(\path{/workspace/hodge-formal-stage1-evidence-20261106/axioms.txt},
source \lean{Stage1Audit.lean}) reports only the standard Lean axioms
\lean{propext}, \lean{Classical.choice}, and \lean{Quot.sound}
(the first two theorems use only \lean{propext} and \lean{Quot.sound}).
The combined Stage~1+Stage~2 audit
(\lean{ResidueAxiomAudit.lean}; log
\lean{residue-audit-e207394.log}) reports
\begin{quote}
\texttt{PASS: 52 incremental theorem declarations; standard axioms only}.
\end{quote}
Repository checks record no \lean{sorry}, \lean{admit}, or
\lean{native_decide} in the incremental modules
(\lean{no-sorry-e207394.log}:
``OK: no sorry/admit/native\_decide, no axiom declarations'').

Documented Stage~1 module build: \path{lake build Hodge.ExplicitResidueAlgebra}
passes with $7084$ build targets
(\path{docs/FULL_FORMALIZATION_STAGE1.md}). The Stage~2 incremental build that
includes this module reports $7085/7085$ success for
\path{Hodge.ExplicitResidueCertificates}
(\lean{build-and-audit-final.log}). A full repository \lean{lake build} at
commit \lean{e207394} completed successfully
(\lean{full-build-e207394.log}).

} % end sloppy

\section{Toolchain pins and reproducibility}
{\sloppy

\begin{center}
\begin{tabular}{@{}ll@{}}
\toprule
Pin & Value \\
\midrule
Lean~\cite{Lean4} & \lean{leanprover/lean4:v4.22.0} (\lean{lean-toolchain}) \\
mathlib input~\cite{mathlib4} & \lean{v4.22.0} (\lean{lakefile.toml}) \\
mathlib checkout & \path{79e94a093aff4a60fb1b1f92d9681e407124c2ca} \\
Git commit & \path{e207394a9ed4255ad804860b8f69fde7e940dd54} \\
PR & \url{https://github.com/BenFrohman/HODGE/pull/36} \\
License & Apache-2.0 (unchanged; no relicense) \\
\bottomrule
\end{tabular}
\end{center}

Reproduction commands (from a checkout of the pinned commit):
\begin{verbatim}
lake exe cache get
lake build Hodge.ExplicitResidueAlgebra
python3 scripts/check_no_sorry.py
# Combined audit after Stage 2 is present:
lake build Hodge.ExplicitResidueAlgebra \
           Hodge.ExplicitResidueCertificates
python3 scripts/check_residue_axioms.py
\end{verbatim}

Local evidence directories inspected while preparing this draft (read-only;
another worker packages proof evidence):
\begin{itemize}[leftmargin=1.4em]
\item \path{/workspace/hodge-formal-stage1-20261106/}
      (mathematical source pin \path{e207394a9ed4255ad804860b8f69fde7e940dd54};
      later inspected worktree \path{fefe2ba9d87c900b538d02c13f7f51f375fcd285}
      retains identical reviewed modules, audit, and toolchain/lock files);
\item \path{/workspace/hodge-formal-stage1-evidence-20261106/}
      (\lean{Stage1Audit.lean}, \lean{axioms.txt},
      \lean{baseline-explicit-curve.lean});
\item \path{/workspace/hodge-formal-stage2-evidence-20261106/}
      and \path{/workspace/hodge-certificate-delivery-20261106/}
      (build/audit logs cited above).
\end{itemize}

} % end sloppy

\paragraph{CI verification at a later identical head.}
User-reported merge of PR~\url{https://github.com/BenFrohman/HODGE/pull/36}
is acknowledged; this packaging records no merge-commit SHA (merge
verification is underway separately). At head
\path{fefe2ba9d87c900b538d02c13f7f51f375fcd285}, where the reviewed Lean
modules, audit, and toolchain/lock match mathematical pin
\path{e207394a9ed4255ad804860b8f69fde7e940dd54}, all five reported GitHub
Actions checks completed successfully:
Lean/build and Lean/attempt
(\url{https://github.com/BenFrohman/HODGE/actions/runs/42317080169}),
Lean certificate/certify on the PR
(\url{https://github.com/BenFrohman/HODGE/actions/runs/42317080616}),
Lean certificate/certify on the push
(\url{https://github.com/BenFrohman/HODGE/actions/runs/42317075046}),
and Paper/paper
(\url{https://github.com/BenFrohman/HODGE/actions/runs/42317080271}).
Both certification logs build the incremental modules and print
\texttt{PASS: 52 incremental theorem declarations; standard axioms only}.
CI success verifies those algebraic declarations; it does not discharge
geometric Hodge/Betti/cycle-class obligations. This note remains a
REVIEW DRAFT with no DOI and no publication claim.

\section{Limitations and relation to the geometric program}

The seven lemmas are genuine polynomial and matrix identities. They do not
instantiate $X$ as a smooth projective scheme, do not produce a Betti cycle
class, and do not identify Jacobian characters with geometric Hodge classes.
Villaflor Loyola's theorem~\cite{Villaflor2021} and Griffiths' Jacobian-ring
comparison remain external classical input
(\lean{NOTICE.md}, \path{docs/FULL_FORMALIZATION_STAGE1.md}).

These algebraic notes can be fully formalized---and, for the seven lemmas
above, already are---while publication of a full geometric Hodge paper for
$X$ remains held pending the missing geometric foundations. No novelty or
peer-review claim is made.

\section*{Author, license, and contact}
Author metadata: \textbf{Frohman, Benjamin Stanley}.
Email: \texttt{frohmanbenjamin@gmail.com}.
ORCID: \href{https://orcid.org/0009-0006-7068-3718}{0009-0006-7068-3718}.
Verified from repository \path{CITATION.cff}, \path{NOTICE}, and
\path{.zenodo.json}.
Copyright (c) 2026 Benjamin Stanley Frohman. Apache License 2.0.
This REVIEW DRAFT assigns no new DOI.

\bibliographystyle{plain}
\bibliography{references}

\end{document}
